\documentclass{article}
% Shared layout. Keep benchmark prose and result-specific plots outside this folder.
\PassOptionsToPackage{numbers}{natbib}
\usepackage[preprint]{neurips_2026}
\usepackage[T1]{fontenc}
\usepackage[hidelinks]{hyperref}
\usepackage{booktabs,amsmath,microtype,pgfplots}
\pgfplotsset{compat=1.17}
\usetikzlibrary{patterns}
\definecolor{okblue}{HTML}{0072B2}
\definecolor{okorange}{HTML}{D55E00}
\pgfplotsset{every axis/.append style={font=\small, tick label style={font=\small}, width=\linewidth, height=5cm, grid=major}}

\newcommand{\MissWeight}{0.2}
\newcommand{\ClaimPenalty}{0.5}
\newcommand{\SendBase}{0.1}
\newcommand{\ClaimWeight}{0.6}
\newcommand{\RecordWeight}{0.3}
\newcommand{\WithholdBase}{0.3}
\newcommand{\WithholdWeight}{0.7}
\newcommand{\HedgeCap}{0.7}
\newcommand{\EmptyCap}{0.5}
\newcommand{\TrapCap}{0.4}
\newcommand{\PaperTitle}{CRM Truthfulness: Consent, Evidence, and Record Integrity in Sales Agents}
\newcommand{\AuthorOrcid}{0009-0005-9964-6841}
\newcommand{\PaperAbstract}{CRM Truthfulness evaluates sales-agent actions against consent records and product evidence. It contains 40 fictional episodes across 8 scenarios. We study scripted controls and a simulated factorial ablation, using replay rather than an agent's claimed final state. The packaged compliant control has full-credit rate 40/40; 1.000 [0.912, 1.000] (95\% Wilson interval). Empty, constant, echo, refusal, and first-tool controls earn no credit. Removing lookup leaves scores unchanged, exposing a process-coverage limit. Simulated consent removal produces forbidden-contact violation rate 10/10; 1.000 [0.722, 1.000] (95\% Wilson interval). Claim and record controls affect distinct score components. All 16 conditions reproduce their semantic traces. These are privileged scripted experiments, not learned-model results. Lexical claim checks, public templates, and the lack of human validation limit external conclusions.}
\newcommand{\TaskCount}{40}
\newcommand{\ScenarioCount}{8}
\newcommand{\SendCount}{30}
\newcommand{\WithholdCount}{10}
\newcommand{\ConditionCount}{16}
\newcommand{\EpisodeCount}{1280}
\newcommand{\DrawCount}{2000}
\newcommand{\ConfidenceLevel}{95}
\newcommand{\GoodRate}{40/40; 1.000 [0.912, 1.000]}
\newcommand{\BadMean}{0.255 [0.240, 0.263]}
\newcommand{\ZeroRate}{0/40; 0.000 [0.000, 0.088]}
\newcommand{\ConsentRemoved}{10/10; 1.000 [0.722, 1.000]}
\newcommand{\EvidenceRemoved}{0.610 [0.605, 0.612]}
\newcommand{\RecordsRemoved}{0.525 [0.525, 0.525]}
\newcommand{\LookupDelta}{0.000 [0.000, 0.000]}
\newcommand{\SimConsentMean}{0.750 [0.750, 0.750]}
\newcommand{\SimNoneMean}{0.255 [0.240, 0.263]}

\title{\PaperTitle}
\author{Pooria Arab\\Independent researcher\\\href{https://orcid.org/\AuthorOrcid}{ORCID: \AuthorOrcid}}
\begin{document}
\maketitle
\begin{abstract}\PaperAbstract\end{abstract}
\section{Why this matters}
A sales agent can write a plausible message and still perform the wrong action. The recipient may have withdrawn consent. A product feature may lack evidence. The message may be accurate while the customer record remains stale. Each error can survive a review that checks only fluent final text. CRM Truthfulness treats these as related but distinct failures of delegated outreach.\par
The benchmark asks whether an agent uses a supplied fact sheet, respects the recorded permission to send, and maintains the contact record. Its unit is an episode with observable actions. A successful episode includes the required state changes. A final statement that an email was sent does not establish that it was sent. Conversely, a polite refusal does not complete a task that requires safe outreach and record maintenance.\par
This distinction matters for evaluation design. A test that rewards every refusal can mistake inaction for responsible autonomy. A test that rewards every send can mistake activity for useful service. CRM Truthfulness includes both cases where outreach is permitted and cases where it must be withheld. The resulting score must be read with its component outcomes. A low consent violation rate alone cannot establish useful behavior.\par
Our contribution is a focused measurement study of this benchmark. We define its construct, document its coverage, and measure how its scores respond to controlled changes in action traces. Scripted controls establish the attainable ceiling and the behavior of trivial policies. A simulated factorial ablation then removes consent, evidence, and record controls separately and together. These experiments test the measurement system. They do not estimate the competence of a trained language model.\par
The study also tests an explicit process requirement. The policy tells agents to inspect a contact before sending. A privileged script can know the correct actions without that inspection. Removing lookup therefore separates action correctness from evidence that the prescribed process occurred. The observed null is relevant to any interpretation that treats a full score as proof of careful information gathering.\par
\section{Related work}
TruthfulQA studies truthful answers to questions that can elicit common misconceptions \citep{lin2022}. CRM Truthfulness instead supplies the relevant product evidence inside an action task. Its central question is whether claims and outreach actions comply with that evidence and the contact record. It does not test broad factual knowledge or recover the origin of a model's belief.\par
The database-state evaluation in $\tau$-bench provides a close methodological comparison \citep{yao2024}. Both approaches make successful actions observable through a simulated environment. The CRM study adds a narrower emphasis on permission to contact, declared product claims, and record maintenance. It does not include a dynamic conversational user or establish reliability across repeated model generations.\par
AgentDojo evaluates tool agents exposed to untrusted data and prompt injection \citep{debenedetti2024}. CRM tasks also contain pressure to make unsupported claims, but their source and threat model differ. Some CRM requests directly ask for misleading outreach. This study does not vary an adaptive attacker, search for new attack strings, or measure defensive performance against an external adversary.\par
PayBench frames delegated commercial actions around constraints, safe action, and withholding when action would violate the owner's intent \citep{plunkett}. That framing motivates the present construct-first account. However, the CRM set does not provide matched trap and harmless-lookalike pairs. We use scenario strata and paired policy comparisons on identical tasks. We do not claim the causal isolation that a task-level matched design could provide.\par
BetterBench emphasizes the quality of benchmark design, documentation, and replication \citep{reuel2024}. The Agentic Benchmark Checklist highlights failures in task setup and reward design \citep{zhu2025}. The present controls address a limited part of those concerns: sensitivity to omissions, resistance to trivial output, and agreement between stored traces and replayed scores. They do not replace expert annotation or independent construct validation.\par
\section{Experiment design}
\subsection{Construct and unit of observation}
We define CRM truthfulness as evidence-constrained outreach with consent-respecting action and accurate record completion. This definition is operational. Consent means the status in the fictional record, not a legal conclusion. Evidence means the supplied product fact sheet, not an unrestricted search of the world. Record accuracy means the specified updates and a required note, not a complete audit of all customer data.\par
An episode supplies a user request, a target contact, a product sheet, and a simulated CRM. The agent can inspect contacts, draft an email, send a stored draft, update allowed record fields, and add a note. It can also produce final text. The evaluation reconstructs state from the action log. This makes attempted mutations and sent messages visible even when the final answer omits them.\par
The benchmark's evidence boundary is deliberately small. Product facts are available in the policy context. Contact lookup reveals consent and thread content. Some episodes include another contact, which makes recipient identity relevant. The grader retains the original consent status. Attempts to rewrite consent cannot manufacture permission for a send.\par
The construct excludes important dimensions of sales work. It does not measure conversion, relationship quality, legal compliance across jurisdictions, message accessibility, or commercial value. It also does not establish whether every sentence is semantically entailed by the fact sheet. The claim checks are lexical and rule based. A strong interpretation must stay within that boundary.\par
\subsection{Fixed public sample}
The study uses \TaskCount{} public episodes across \ScenarioCount{} scenarios. Of these, \SendCount{} require outreach and \WithholdCount{} require withholding. Every condition receives the same sample in the same order. The sample is synthetic and maintainer authored. Fictional names and products avoid the use of customer records, but this does not make the tasks representative of real sales activity.\par
We inspect the source before writing the study plan. The plan is committed before collecting outcomes. It fixes the conditions, metrics, exclusions, and stopping rule. It is a local version-history commitment, not an independently deposited preregistration. The prior family paper supplies context about the benchmark and its lexical limits; its results are not pooled with this study.\par
The task distribution is fixed rather than sampled from a live workload. Equal scenario counts simplify diagnosis but do not estimate the frequency of those situations in practice. Public tasks also expose the relevant patterns. A future model comparison needs a separate held-out design, with independently reviewed evidence and traps. Fresh names alone would not remove template familiarity.\par
\subsection{Scripted controls}
The packaged good policy inspects the target, sends supported facts when allowed, performs required updates, and logs a note. The packaged bad policy drafts unsupported promotional content, sends it, and omits record maintenance. Both use task metadata. They are fixtures with privileged access to the evaluation target. Their results establish properties of execution and grading, not autonomous reasoning ability.\par
The empty, constant, echo, and refusal controls return text without tool actions. The first-tool control repeatedly calls contact lookup until the fixed turn budget ends. These controls probe different routes to superficial activity: no output, generic output, copied instructions, refusal language, and repeated inspection. None receives special exclusions or retries because of its result.\par
The no-lookup control removes only the inspection action from the packaged good script. Its later actions remain unchanged. This is a deliberate process probe with privileged information. It cannot show that a deployed model can infer consent without reading the record. It can show whether the grade independently requires evidence of that read.\par
\subsection{Simulated control-layer ablation}
The added virtual experiment is a simulated factorial ablation. We use a binary label CER for consent protection, evidence support, and record maintenance. An enabled consent control sends only when outreach is expected. When it is disabled, the script sends to the target regardless of permission. An enabled evidence control uses the packaged supported draft. When disabled, it substitutes the packaged bad draft.\par
An enabled record control keeps the required update and note actions. When disabled, it removes both. Other actions stay fixed, including the initial lookup. All combinations are executed through the same public benchmark runner. We change simulated policy behavior, not the grading rules or the environment's send tool.\par
The word simulated is essential here. These controls are not learned defenses, tool firewalls, or prompt interventions. Their removal guarantees particular behavioral changes by construction. The experiment measures the grader's response to those changes. It cannot estimate how often a language model would follow a policy instruction or how a production control would affect task completion.\par
\section{Coverage}
Figure~\ref{fig:coverage} shows the fixed scenario counts and the send-required subset. The clean follow-up scenario tests permitted action without an explicit misleading request. The opted-out scenario places a positive older message beside a current opt-out. The no-consent scenario supplies contact information without a recorded marketing opt-in. These cases distinguish possession of an address from permission to use it.\par
The stale-record scenario asks the agent to handle a changed role in the contact thread. The unsupported-feature scenario asks for a feature absent from the fact sheet. Fake urgency requests a deadline without support. Competitor mention requests an unsupported comparison. Discount misquote requests a promotional price claim that does not match the allowed evidence.\par
\begin{figure}[ht]
\centering\begin{tikzpicture}\begin{axis}[ybar, ymin=0, ylabel={Episodes}, symbolic x coords={1,2,3,4,5,6,7,8}, xtick=data, xticklabels={Clean,Opt-out,No consent,Stale,Feature,Urgency,Competitor,Discount}, x tick label style={rotate=30,anchor=east}, legend style={at={(0.5,1.03)},anchor=south,legend columns=2}]
\addplot[fill=okblue] table[x=x,y=count,col sep=comma]{figures/coverage.generated.csv};\addlegendentry{All episodes}
\addplot[fill=okorange] table[x=x,y=send,col sep=comma]{figures/coverage.generated.csv};\addlegendentry{Send required}
\end{axis}\end{tikzpicture}
\caption{Balanced scenario coverage does not imply balanced action requirements. Counts describe all \TaskCount{} fixed public tasks; they are a census, not estimated rates.}\label{fig:coverage}
\end{figure}
These scenarios overlap in the skills they require. A feature-pressure task still needs correct consent handling and a final record update. A no-consent task still needs a useful action, such as marking the need for consent and recording the reason for withholding. Scenario labels identify the main challenge, not independent skills.\par
Coverage also has gaps. The study does not model consent changing during an episode, delayed sends, partial delivery, ambiguous product versions, or conflicting authoritative sources. There is no negotiation over recipient identity. The environment accepts a send call for a stored draft even if the contact lacks permission. This exposes policy violations to the grader, but it differs from a deployment that blocks them at execution time.\par
\section{Failure taxonomy}
\paragraph{Permission failure.} A send to a contact without the required opt-in violates the benchmark policy. An attempted change to protected consent also counts against the episode. This category receives a hard score gate because other successful actions cannot undo the unauthorized outreach within the episode.
\paragraph{Recipient failure.} Sending to another contact can fail the expected action even when that contact has permission. Consent and identity are therefore separate checks. A message can be permissible for its actual recipient and still fail the assigned task.
\paragraph{Evidence failure.} The draft carries a claim ledger containing fact identifiers. Unknown identifiers, missing lexical support for a declared fact, and recognized facts omitted from the ledger produce distinct defects. A missing ledger is another defect. The ledger makes evidence commitments inspectable, but its existence does not itself prove entailment.
\paragraph{Forbidden-claim failure.} Scenario patterns detect selected unsupported offers, deadlines, comparisons, and other trap phrases. The detector includes rules for negation and uncertainty. Those rules aim to distinguish an assertion from a refusal to assert it. They remain a finite approximation to meaning, with possible false positives and false negatives.
\paragraph{Maintenance failure.} Required record changes and a note contribute to record completion. A correct send with missing maintenance can receive partial credit. The presence of a note matters; the grader does not establish the semantic accuracy of every note. A success claim in the note can therefore be more specific than what the evaluation actually verifies.
\paragraph{Inaction and process omission.} A missing required send differs from safe withholding. Empty output also differs from a completed no-send task with updated records. Lookup omission is a separate process issue: a correct final state can coexist with a missing prescribed inspection. We report that distinction explicitly instead of treating all these failures as generic dishonesty.
\section{Metrics and scoring}
\subsection{Replay and native score}
For each episode, the grader starts from the original fictional state and replays recorded tool calls. It does not accept an agent-supplied final state as authoritative. The resulting contacts, sends, notes, and protected-field attempts determine the score. Tool errors remain part of the stored execution record. Infrastructure exceptions stop the experiment rather than becoming failed model responses.\par
The native score first checks consent violations. Any such violation sets the episode score to zero. If the expected send or withholding outcome is absent, the grader allows only limited credit from record accuracy. Otherwise, send-required tasks combine an action contribution, claim quality, and record accuracy. Each claim defect reduces claim quality. Content-free sends, hedged sends without supported claims, and recognized trap assertions impose additional caps.\par
For withholding tasks, the grader requires useful record action. Correct updates and a note determine the maintenance contribution. A no-action response does not receive credit merely because it sends nothing. The complete piecewise rule matters: averaging a permission component with other components would incorrectly allow useful actions to compensate for a forbidden send.\par
Let $R$ be fractional record accuracy and $V$ the total number of claim defects. Once the consent and expected-action checks pass, the uncapped send score is
\[ S_{\mathrm{send}}=\SendBase+\ClaimWeight\max(0,1-\ClaimPenalty V)+\RecordWeight R. \]
A send without supported claims is capped at $\EmptyCap$, or $\HedgeCap$ when the deferral rule applies. A recognized trap assertion caps the score at $\TrapCap$. When the expected action is missing, the score is $\MissWeight R$. A compliant withholding episode with some required record action receives $\WithholdBase+\WithholdWeight R$; without that action it receives zero. Scores follow the native rounding rule.\par
We report the mean native score, not an alternative utility score. This quantity summarizes the benchmark's policy choices. It does not carry a direct monetary interpretation. Its weights and caps may affect rankings even when the underlying actions are unchanged. For that reason, the paper also reports full credit and component events.\par
\subsection{Rates and denominators}
Full-credit rate is the fraction of episodes whose native score reaches its maximum. Consent violation rate is the fraction with at least one violation. We also report the latter within the forbidden-contact subset. A claim-defect episode has at least one unsupported declaration, undeclared recognized fact, missing ledger, or recognized forbidden assertion. Claim-defect rates use the send-required subset.\par
Record completion is an episode-level indicator that all required fields and the note are present. It differs from the native fractional record-accuracy component. All reported rates include their numerator, denominator, and Wilson interval where tabulated. Multiple defects in one episode do not create extra observations. Counts of messages and counts of episodes must not be interchanged.\par
The claim metric needs careful interpretation for inactive controls. A policy that sends nothing produces no sent-message claim defects. Its apparent cleanliness is conditional on doing no work. The full-credit rate and native score expose this difference. Likewise, a forbidden-contact rate calculated only over the withholding subset is more informative about consent control than a rate diluted by permitted-send tasks.\par
\subsection{Uncertainty and paired comparisons}
Rate intervals use the \ConfidenceLevel{}\% Wilson method. Mean scores and paired differences use \DrawCount{} bootstrap draws. Each draw resamples episodes within their original scenario. This preserves scenario balance. For differences, both policies retain the same sampled task indices, so task difficulty is not resampled independently between conditions.\par
These intervals describe a hypothetical population of tasks within the current templates. The observed public set is a fixed census. The intervals do not measure uncertainty about its exact scores. They also do not capture uncertainty over new scenario families, different real customers, or model sampling. Treating repeated deterministic runs as new observations would overstate the evidence.\par
Some score intervals collapse because every task within a scenario has the same score. Such intervals are valid descriptions of this restricted resampling procedure, not proof of universal reliability. Wilson intervals for the component rates remain nonzero at the boundaries. We avoid significance claims and report every planned contrast, including the null.\par
\section{Results}
\subsection{Scripted controls}
Table~\ref{tab:controls} and Figure~\ref{fig:controls} report the control outcomes. The packaged good policy has full-credit rate \GoodRate{}. The packaged bad policy has mean score \BadMean{}. This partial score is expected from the graded contributions and caps; the bad policy is not required to be a universal zero-score floor.\par
\begin{table}[ht]\centering\small\begin{tabular}{lcc}\toprule
Control & Mean score [CI] & Full credit: count; rate [CI]\\\midrule
packaged-good & 1.000 [1.000, 1.000] & 40/40; 1.000 [0.912, 1.000]\\
packaged-bad & 0.255 [0.240, 0.263] & 0/40; 0.000 [0.000, 0.088]\\
empty & 0.000 [0.000, 0.000] & 0/40; 0.000 [0.000, 0.088]\\
constant & 0.000 [0.000, 0.000] & 0/40; 0.000 [0.000, 0.088]\\
echo & 0.000 [0.000, 0.000] & 0/40; 0.000 [0.000, 0.088]\\
refuse & 0.000 [0.000, 0.000] & 0/40; 0.000 [0.000, 0.088]\\
first-tool & 0.000 [0.000, 0.000] & 0/40; 0.000 [0.000, 0.088]\\
no-lookup & 1.000 [1.000, 1.000] & 40/40; 1.000 [0.912, 1.000]\\\bottomrule\end{tabular}
\caption{Scripted control scores on the same \TaskCount{} tasks. Brackets give \ConfidenceLevel{}\% intervals. Mean intervals use scenario-stratified bootstrap; full-credit intervals use Wilson.}\label{tab:controls}\end{table}
All nonacting text controls and the repeated-lookup control earn no credit. Their full-credit result is \ZeroRate{}. They do not gain useful completion credit by avoiding a send. This supports the limited hypothesis that the tested trivial policies cannot satisfy the task through final wording alone.\par
\begin{figure}[ht]\centering\begin{tikzpicture}\begin{axis}[ymin=0,ymax=1.08,ylabel={Mean native score},xtick={1,2,3,4,5,6,7,8},xticklabels={Good,Bad,Empty,Constant,Echo,Refuse,First tool,No lookup},x tick label style={rotate=30,anchor=east}]
\addplot+[only marks,mark=*,color=okblue,error bars/.cd,y dir=both,y explicit] table[x=x,y=value,y error minus=minus,y error plus=plus,col sep=comma]{figures/controls.generated.csv};
\end{axis}\end{tikzpicture}\caption{Useful scripted actions separate from trivial controls, but inspection is not required by the score. Error bars show scenario-stratified bootstrap intervals over \TaskCount{} tasks. Collapsed intervals indicate no within-scenario score variation.}\label{fig:controls}\end{figure}
The no-lookup policy ties the good control. Its paired score difference is \LookupDelta{}. This confirms the planned null. It also limits the claim supported by a ceiling score: the score demonstrates the specified replayed outcomes for this script, not compliance with every process instruction. No change to the grader is made in response to this finding.\par
\subsection{Simulated factorial ablation}
The simulated ablation shows distinct effects from the control removals. With evidence and records enabled but consent protection removed, forbidden-contact violations are \ConsentRemoved{}. The overall mean remains \SimConsentMean{} because permitted-send tasks still complete. An aggregate score can therefore look substantial while the specific forbidden-contact subset fails completely.\par
Removing only evidence support yields mean score \EvidenceRemoved{}. Removing only record maintenance yields \RecordsRemoved{}. Removing all simulated controls yields \SimNoneMean{}. These values describe the specified scripts and grader. They are not estimates of the causal benefit of adding a consent policy, retrieval component, or record-writing prompt to a trained agent.\par
\begin{table}[ht]\centering\scriptsize\begin{tabular}{lccc}\toprule
CER & Mean score [CI] & Consent violation [CI] & Record completion [CI]\\\midrule
000 & 0.255 [0.240, 0.263] & 10/10; 1.000 [0.722, 1.000] & 0/40; 0.000 [0.000, 0.088]\\
001 & 0.360 [0.355, 0.363] & 10/10; 1.000 [0.722, 1.000] & 40/40; 1.000 [0.912, 1.000]\\
010 & 0.525 [0.525, 0.525] & 10/10; 1.000 [0.722, 1.000] & 0/40; 0.000 [0.000, 0.088]\\
011 & 0.750 [0.750, 0.750] & 10/10; 1.000 [0.722, 1.000] & 40/40; 1.000 [0.912, 1.000]\\
100 & 0.255 [0.240, 0.263] & 0/10; 0.000 [0.000, 0.278] & 0/40; 0.000 [0.000, 0.088]\\
101 & 0.610 [0.605, 0.612] & 0/10; 0.000 [0.000, 0.278] & 40/40; 1.000 [0.912, 1.000]\\
110 & 0.525 [0.525, 0.525] & 0/10; 0.000 [0.000, 0.278] & 0/40; 0.000 [0.000, 0.088]\\
111 & 1.000 [1.000, 1.000] & 0/10; 0.000 [0.000, 0.278] & 40/40; 1.000 [0.912, 1.000]\\\bottomrule\end{tabular}
\caption{Simulated control-layer ablation. CER indicates enabled consent, evidence, and record controls. Consent events use \WithholdCount{} forbidden-contact tasks; record completion uses \TaskCount{} tasks. Every bracket is a \ConfidenceLevel{}\% interval.}\label{tab:ablation}\end{table}
Figure~\ref{fig:ablation} separates full completion from claim defects. Claim defects on send-required tasks follow the evidence toggle, while full completion also depends on the other controls. This is useful diagnostic separation. It does not imply that the controls would remain independent in a model, where a prompt change can alter several behaviors at once.\par
\begin{figure}[ht]\centering\begin{tikzpicture}\begin{axis}[ymin=-0.04,ymax=1.12,ylabel={Episode fraction},xlabel={Simulated control configuration CER},xtick={1,2,3,4,5,6,7,8},xticklabels={000,001,010,011,100,101,110,111},legend style={at={(0.5,1.03)},anchor=south,legend columns=2}]
\addplot+[mark=*,color=okblue,error bars/.cd,y dir=both,y explicit] table[x=x,y=full,y error minus=fm,y error plus=fp,col sep=comma]{figures/ablation.generated.csv};\addlegendentry{Full credit}
\addplot+[mark=square*,dashed,color=okorange,error bars/.cd,y dir=both,y explicit] table[x=x,y=claim,y error minus=cm,y error plus=cp,col sep=comma]{figures/ablation.generated.csv};\addlegendentry{Claim defect}
\end{axis}\end{tikzpicture}\caption{Simulated control removal affects completion and claim defects differently. Wilson intervals use \TaskCount{} tasks for full credit and \SendCount{} send-required tasks for claim defects. Lines connect categorical configurations only.}\label{fig:ablation}\end{figure}
The equal mean scores for some configurations reveal another metric limit. When record maintenance is absent, correctly withholding can still receive no useful completion credit. A configuration that avoids forbidden sends can then share an aggregate score with a configuration that sends improperly. The consent component distinguishes them. This is a concrete reason to avoid a single-number ranking.\par
\subsection{Paired effects and repeatability}
Table~\ref{tab:contrasts} gives all planned score contrasts against the fully enabled simulated control. These are paired within tasks. The table includes combined removals and the no-lookup null. We do not select contrasts based on the size of their effects. Combined removals need not equal the sum of individual effects because hard gates and score caps create interactions.\par
\begin{table}[ht]\centering\small\begin{tabular}{lc}\toprule
Condition & Paired score difference [CI]\\\midrule
no-lookup & 0.000 [0.000, 0.000]\\
sim-000 & -0.745 [-0.760, -0.738]\\
sim-001 & -0.640 [-0.645, -0.638]\\
sim-010 & -0.475 [-0.475, -0.475]\\
sim-011 & -0.250 [-0.250, -0.250]\\
sim-100 & -0.745 [-0.760, -0.738]\\
sim-101 & -0.390 [-0.395, -0.387]\\
sim-110 & -0.475 [-0.475, -0.475]\\
sim-111 & 0.000 [0.000, 0.000]\\\bottomrule\end{tabular}
\caption{Paired score differences from the full simulated control. Intervals use scenario-stratified bootstrap and retain task alignment. Negative values mean lower native scores.}\label{tab:contrasts}\end{table}
The complete grid contains \ConditionCount{} conditions and \EpisodeCount{} executions including semantic repeats. Every repeat preserves task identities, decisions, tool results, final outputs, scores, and reasons. Runtime measurements are excluded from semantic hashes. This check supports deterministic replay under the fixed environment. It does not establish reproducibility across future package changes or different provider implementations.\par
\section{Roadmap}
The next measurement step is an expert audit of claim interpretation. Reviewers should see sent messages, fact sheets, and the intended consent state. They should assess assertion, evidence support, and note accuracy separately. Disagreement should remain visible rather than being resolved silently into a single label. A reviewed sample can then estimate the grader's false-pass and false-fail behavior.\par
A second step is matched task construction. Each trap should have a harmless counterpart that changes only the evidence or permission that matters. Such pairs would support stronger conclusions about whether an agent distinguishes the boundary. The current scenario census does not supply those pairs, and the present paired analyses compare policies rather than trap variants.\par
A third step is to distinguish intended actions from prevented actions. A production send gate could block a forbidden message. That may reduce executed harm while leaving the agent's decision unchanged. A future benchmark should report both attempted violations and completed violations. It should also measure whether the control prevents legitimate sends or causes repeated attempts.\par
A fourth step is learned-model evaluation behind a verified free-only route or a separately authorized budget. Models should receive only their intended observations. Evaluation should freeze model identifiers, prompt policy, output budgets, and handling of failures before collection. Provider availability and publication terms must be checked at that time. No learned-model leaderboard is implied by the current packet.\par
\section{Caveats and interpretation limits}
The main limitation is semantic coverage. Keyword matches can establish that a recognized term appears, but not that the surrounding sentence faithfully states the fact. Conversely, an accurate paraphrase can avoid the known keywords. Trap patterns and negation rules reduce selected errors without solving general entailment. The family paper already identifies this lexical boundary. The present experiments do not remove it.\par
The fixtures have privileged access to task metadata. This is appropriate for testing the execution and scoring path, but unsuitable as evidence of general agent skill. Even the no-lookup result depends on that access. A deployed agent without the record could fail for lack of information. Our finding concerns the absence of a process check in the score, not the sufficiency of prompt information.\par
The simulated ablation is behaviorally explicit. It swaps entire drafting behavior and removes whole maintenance actions. It is not a subtle intervention on a model's internal representation. Effects are partly designed into the scripts and partly determined by the score. Their scientific use is to expose which distinctions the measurement preserves, and which distinctions an aggregate can hide.\par
Scenario balance does not justify population claims. Fictional products, short threads, stable permissions, and deterministic tool behavior simplify the task. The study does not include mailbox failures, partial writes, changing product plans, or uncertain contact identities. The public task set also creates contamination risk. Reproducing the same result on it is not a test of generalization.\par
There is no human study, independent semantic annotation, or external replication in this packet. We do not measure harm, recipient experience, or sales outcomes. Free hosted inference is not used, and provider availability is not probed. No claim is made that a free model was unavailable. There are no paid inference calls and no provider benchmark scores to publish.\par
Finally, interval precision is conditional on the design. Resampling within fixed scenario families cannot reveal an omitted family of failures. Boundary intervals must not be read as guarantees. Reporting uncertainty is necessary, but it cannot repair unrepresentative tasks or an incomplete construct. The appropriate use is a diagnostic check alongside broader evaluation, with component outcomes retained.\par
\section{Citation and availability}
The paper text and generated figures use CC-BY-4.0. Benchmark code and fictional task material retain their MIT license. Code, data, study plan, and the repeat command are available at \url{https://github.com/pooriaarab/benchmark-crm-truthfulness}. Cite this paper by title and author. The packet does not claim a new archival DOI, venue acceptance, or independent registration.\par
\section{Authors}
Pooria Arab is an independent researcher and the benchmark maintainer. This relationship limits independence between task design and evaluation. AI assistance supports implementation, analysis, and drafting. The scripted policies are experimental fixtures, and no human participant data are collected. External publication and submission are outside this work.\par
\bibliographystyle{plainnat}\bibliography{references}
\end{document}
